\chapter{Glossary of XNA and CNA Terms}
\label{ch:glossary}

\begin{description}
\item[BGFX] One of CNA's eleven graphics backends; a third-party rendering library integrated
  directly via its native API. See Volume~I, Chapter~20.
\item[CANVAS] CNA's Emscripten-only, HTML \texttt{<canvas>}-based 2D graphics backend. See
  Volume~I, Chapter~23.
\item[CNA] The project this book is about: a C\texttt{++}23 reimplementation of the XNA~4.0
  programming model, built on SDL3 with a pluggable graphics-backend layer.
\item[DX3] CNA's graphics backend fronting the sibling \texttt{free-direct} project; a 2D-only
  backend, in the same spirit as \texttt{SDL\_RENDERER}. See Volume~I, Chapter~24.
\item[EASYGL] CNA's OpenGL/OpenGLES graphics backend, built on the sibling \texttt{easy-gl}
  project; CNA's most feature-complete backend by its own account. See Volume~I, Chapter~18.
\item[Effect] The XNA base class for a shader program bound to a \cnaclass{GraphicsDevice}.
  CNA ships five stock effects and, on three backends, a runtime-compiled custom
  \texttt{ShaderEffect} path --- see the shader/.fx gap chapter, Volume~I, Chapter~15.
\item[FNA] An open-source, C\#, SDL2-based reimplementation of the XNA~4.0 API, maintained
  independently of CNA. FNA is CNA's primary behavioral reference and the subject of CNA's
  own differential-testing infrastructure --- not the same thing as the real Microsoft
  XNA~4.0 specification, a distinction this book returns to repeatedly (most sharply in the
  Media chapter's \cnaclass{Song} finding, Chapter~3 of this volume).
\item[GamerServices] The XNA namespace covering Xbox-Live-era player profile, achievement,
  leaderboard, and system-UI concepts. CNA implements it with real local persistence where
  FNA does not implement it at all. See Chapter~5 of this volume.
\item[GraphicsDevice] The central class every rendering operation in CNA passes through;
  covered in Volume~I, Chapter~9.
\item[.mdxb / .xnb] Real XNA's compiled binary content-pipeline output format. CNA does not
  implement a general \texttt{.xnb} reader by design (with one real exception, the
  \cnaclass{Model} reader) --- see Volume~I, Chapter~8, and this volume's migration guide,
  Chapter~17.
\item[MojoShader] The bytecode-to-GLSL cross-compiler FNA uses to run compiled Direct3D~9
  shader bytecode. CNA has no equivalent, which is the root cause of the shader/.fx bytecode
  gap covered in Volume~I, Chapter~15.
\item[NOXNA] CNA's marker convention (a macro, and a matching compile-time build option) for
  any declaration inside the XNA namespace tree that is not part of the real XNA~4.0 API. See
  Volume~I, Chapter~3.
\item[Oracle (D3D9)] The pixel-for-pixel verification methodology comparing CNA's D3D9
  backend against a genuinely running XNA~4.0 runtime under Wine, at zero tolerance. See
  Volume~I, Chapter~22, and this volume's Windows/Wine chapter, Chapter~13.
\item[SDL3] The cross-platform windowing/input/audio library CNA is built on. Not itself a
  graphics backend --- \texttt{SDL\_RENDERER} (one of CNA's eleven backends) is built on
  SDL3's own 2D renderer specifically.
\item[SharpRuntime / sharp-runtime] The sibling C\texttt{++} reimplementation of a subset of
  the .NET base class library that CNA depends on for its non-XNA-specific vocabulary
  (\texttt{System::Object}, \texttt{TimeSpan}, \texttt{EventHandler<T>}, and more). Covered in
  full in Part~II of this volume.
\item[SkinnedModelEXT] The CNA-original, \texttt{NOXNA} real-rendering extension for the
  Avatar system, structurally separate from the faithful, inert base Avatar API. See
  Chapter~7 of this volume.
\item[SpriteBatch] The primary 2D rendering API in XNA and CNA alike. See Volume~I, Chapter~10.
\item[Stock effects] The five built-in XNA shader effects --- \cnaclass{BasicEffect},
  \cnaclass{AlphaTestEffect}, \cnaclass{DualTextureEffect}, \cnaclass{EnvironmentMapEffect},
  \cnaclass{SkinnedEffect}. See Volume~I, Chapter~13.
\item[VULKAN] One of CNA's eleven graphics backends, targeting the Vulkan API directly. See
  Volume~I, Chapter~19.
\item[WEBGPU] CNA's experimental fifth graphics backend, targeting native \texttt{wgpu-native}.
  See Volume~I, Chapter~21.
\item[XACT] Microsoft's Cross-platform Audio Creation Tool format/engine, the format CNA's own
  Audio namespace parses and plays back through SDL3\_mixer rather than a FAudio-equivalent
  engine. See Chapter~2 of this volume.
\item[xna4-spec] The sibling repository providing a complete, machine-readable XML
  transcription of the real XNA~4.0 API documentation, used to audit CNA's coverage against a
  ground truth rather than against FNA's own source. See Chapter~19 of this volume.
\end{description}
